% Legacy standalone literature-review chapter kept only for reference.
% The active thesis structure has merged this content into Chapter 1 (绪论)
% to align with the本科毕业设计论文范文.docx chapter format.
\chapter{国内外研究现状}

\section{国外研究现状}

国际上围绕雾天气视觉感知的研究大体可以分为四条路线：基于图像恢复的去雾路线、基于能见度或雾参数估计的物理路线、面向恶劣天气鲁棒感知的目标检测路线，以及强调天气自适应表示学习和多天气统一建模的综合路线。

在图像恢复与去雾研究方面，早期方法主要建立在暗通道先验、颜色衰减先验、对比度增强等统计规律基础上，能够在一定程度上恢复图像清晰度，但对真实交通场景中的局部强光、运动模糊、压缩噪声和空间非均匀雾分布适应性有限。进入深度学习阶段后，研究重点逐渐转向端到端的图像去雾网络与恢复-检测一体化方法。2025 年的综述论文 Clear Roads, Clear Vision: Advancements in Multi-Weather Restoration for Smart Transportation 对交通场景下多天气恢复技术进行了系统总结，指出雾、雨、雪、低光照等视觉退化并非彼此孤立，而是共同影响智能交通系统感知稳定性的关键因素。该文的价值在于从系统视角总结了恢复模型与下游感知任务之间的关系，说明仅以主观清晰度作为目标并不足以支撑道路安全应用。

在雾参数与能见度估计方面，国外研究越来越重视“由图像恢复转向参数估计”的趋势。2025 年发表于 arXiv 和后续数据库收录的 Estimating Fog Parameters from a Sequence of Stereo Images 提出从立体图像序列中联合估计雾参数的优化框架，强调同时求解雾相关参数可以降低误差传播。与传统逐步估计方式相比，这类方法更加关注雾过程本身的物理可解释性，对本文利用 beta 作为回归目标具有直接参考意义。尽管该工作采用的是双目序列场景，而本文依赖单目监控视频与单目深度估计，但其核心思想说明：若能够从视觉退化中提取稳定参数，就可以将模糊的“有雾/无雾”判断转化为更具工程价值的连续量化指标。

在雾天目标检测方面，近两年国际研究明显向“去雾与检测融合”方向发展。Oxford Academic 发表的 Defog YOLO for road object detection in foggy weather、Pattern Recognition 上的 Dehazing \& Reasoning YOLO、arXiv 上的 FogGuard、D-YOLO 与 ADAM-Dehaze 等方法均试图把图像恢复或雾先验嵌入到检测网络中，而不是把去雾作为完全独立的前置步骤。这类方法的共同特点是：一方面利用大气散射模型、先验知识或可学习的恢复模块提升图像特征质量；另一方面直接以检测性能为优化目标，避免出现“图像看起来更清晰，但检测结果并未改善”的问题。2026 年发表的 Fog-YOLO: A Real-Time Object Detector for Foggy Autonomous Driving Scenarios 进一步表明，实时性仍然是雾天检测算法的重要约束，尤其在道路视频监测和边缘计算场景中，模型复杂度不能无限增加。

与此同时，国外学界还在推动更广义的恶劣天气鲁棒感知研究。诸如 From Fog to Failure、Robust ADAS、WARLearn 等工作把雾与雨、雪、低光照共同纳入模型鲁棒性分析框架中，强调跨天气域迁移、天气自适应表征学习和下游任务性能保障。这类研究为本文提供了两点启发：其一，天气识别不应只是孤立分类任务，而应能够反过来调节检测策略；其二，单一晴天数据训练的检测器在雾场景下存在明显域偏移，因此在训练阶段引入天气合成与联合优化是必要的。

总体来看，国外研究在恢复模型、恶劣天气检测与参数估计方面发展迅速，但仍存在几个不足：第一，不少方法依赖自动驾驶前视相机数据，与固定监控视频存在视角差异；第二，大量工作集中于恢复后检测而非直接面向交通业务指标；第三，许多方法仅输出检测结果而缺乏可解释的雾强度量化；第四，面向固定交通监控系统的统一多任务建模仍相对不足。本文的工作正是在这一背景下，尝试将检测、天气识别和雾浓度估计在同一监控视频场景中统一起来。

\section{国内研究现状}

国内关于高速公路低能见度与团雾监测的研究具有明显的工程应用导向。一方面，国内高速公路管理部门、气象部门和公安交管部门围绕恶劣天气高影响路段持续推进标准建设与联动机制；另一方面，学术界和产业界围绕雾天识别、交通气象预警、图像去雾与低能见度目标检测也开展了大量研究。

从行业标准和业务体系建设来看，近两年我国在交通气象服务方面进展明显。中国气象局于 2024 年发布行业标准 QX/T 729—2024《高速公路交通安全管控天气风险预警等级》，对高速公路交通安全管控中的天气风险预警等级进行了规范；2026 年初又正式发布 QX/T 808—2025《高速公路雾天视频图像能见度等级识别技术规范》，明确了雾天视频图像数据采集、识别模型构建和结果验证等技术要求。结合中国气象局与交通运输、公安交管等部门联合推进的“一路多方”机制，以及 2025 年试点高速公路浓雾和道路结冰风险地图发布等工作，可以看出我国正在逐步形成从气象观测、标准制定到路段风险识别和管控处置的完整业务链。对于本文而言，这意味着团雾监测系统不仅具有算法研究价值，也与当前行业标准和管理需求高度契合。

在国内论文研究方面，近两年相关成果主要集中于三个方向。

第一，交通气象阈值和预警机制研究。2024 年《雾天条件下交通气象预警阈值探究》围绕雾天交通运行中的预警阈值进行了分析，说明在低能见度条件下，道路管理更需要分级、分场景的风险表达方式，而不是单一固定阈值。这一点与本文将雾状态同时表示为离散类别和连续 beta 的做法在思路上是一致的：既要便于业务决策，也要保留连续变化信息。

第二，面向交通场景的 AI 事件检测与智能感知研究。2024 年《大模型在高速公路事件检测中的应用探索》表明，国内研究和工程界已经开始将多模态和大模型思路引入高速公路视频事件分析。虽然该类研究未必直接聚焦团雾，但它体现了当前国内高速公路智能监测正在从单一目标识别向复杂场景理解演进。本文的统一多任务方案可以被视为这种“从单点任务到场景级感知”趋势下的一种轻量化、可部署化实现路径。

第三，雾天图像去雾和跨域目标检测研究。2025 年《基于 CycleGAN 的无监督图像去雾网络》《结合特征增强注意力的混合卷积去雾网络》以及 UMT-HPLA: 半监督雾天跨域目标检测模型等工作表明，国内学者在无监督去雾、特征增强、跨域鲁棒检测方面持续发力。这类研究多数尝试解决“合成雾与真实雾分布不一致”“晴天训练模型在雾天泛化差”“去雾网络与检测网络割裂”等问题，与本文利用在线造雾和多任务联合训练缓解域偏移的方法具有较强关联。

此外，从管理实践上看，安徽、江苏、福建、河南、湖北、四川等地在 2024—2025 年间持续推进高速公路恶劣天气监测、分级预警和联动管控。交通运输部和中国气象局公开资料显示，部分省份已经形成团雾多发路段更高密度的交通气象站网，并开始探索气象风险地图、数字孪生和多部门联动处置。可以说，国内在“业务需求牵引算法研究”这一点上具有明显优势。

尽管如此，国内研究仍存在一些短板。其一，公开、规范、可复现的高速公路团雾视频数据集仍然不足，导致许多方法难以横向对比；其二，很多方法关注去雾质量，却没有把雾分类、浓度估计和目标检测整合为统一输出；其三，部分模型面向论文指标设计，工程链路不完整，难以直接迁移至监控中心现有视频系统；其四，真实业务更需要“雾状态 + 车辆状态 + 风险提示”的组合输出，而非单一去雾图像。

\section{现有研究存在的主要问题}

综合国内外研究现状，可以发现现有工作虽然已经取得大量成果，但距离“固定交通监控视频中的可部署团雾监测系统”仍有明显差距。

第一，数据问题仍是核心瓶颈。真实团雾数据稀缺、标签不统一、采集成本高，导致很多研究依赖合成雾或通用自动驾驶数据，缺少与高速公路固定监控场景高度匹配的数据支撑。

第二，任务组织方式仍较为分散。已有工作往往分别研究去雾、检测、分类或能见度估计，但真正面向业务时，管理人员需要的是统一、联动、可解释的结果输出。

第三，工程化程度不足。许多论文方法没有公开完整代码，或者只给出模型结构而缺乏训练、验证、推理、导出与部署闭环，不利于毕业设计和工程实践验证。

第四，缺少对“可解释雾强度指标”的重视。实际交通管控更需要“当前雾有多重”“风险是否在上升”“是否应当调低检测阈值或触发诱导信息”等连续判断，因此单纯二分类或三分类结果并不充分。

正是基于上述不足，本文提出采用深度估计辅助的在线造雾策略与统一多任务模型，在不引入复杂多传感器前提下尽量提升固定监控视频场景中团雾监测的完整性与可部署性。
